% !TEX root = ../../../main.tex
\section{Characteristic polynomials and similarity}
\label{sec:linear-algebra-i:characteristic-polynomials-and-similarity}

% Sources: IMG_9269, IMG_9320--IMG_9322;
% October 15 PDF, pp. 6--10 and matrix example p. 26.
\subsection{Characteristic polynomials}

\begin{definition}
  The characteristic polynomial of \(M\in\mathbb F^{n\times n}\) is
  \[
    p_M(x)=\det(M-xI).
  \]
  For an endomorphism \(L:V\to V\), its characteristic polynomial is
  \(p_L(x)=\det(L-xI)\).
\end{definition}

Let \(\delta^i_j\) denote the entry in row \(i\), column \(j\) of
the identity matrix,
\[
  \delta^i_j=\begin{cases}1&\text{if }i=j,\\0&\text{if }i\ne j.\end{cases}
\]
The permutation formula gives
\[
  p_M(x)=\sum_{\sigma\in S_n}\varepsilon(\sigma)
    \prod_{j=1}^n\bigl(M^{\sigma(j)}_j-x\delta^{\sigma(j)}_j\bigr).
\]
This is a polynomial of degree \(n\). Its leading term is \((-x)^n\),
and its constant term is \(\det M\). Its roots in the scalar field
are exactly the eigenvalues of \(M\).

There is also a direct expansion for \(L\). Normalize a top-degree form
\(D\) on a basis \((e_1,\ldots,e_n)\). Then
\begin{align*}
  p_L(x)
    &=D((L-xI)e_1,\ldots,(L-xI)e_n)\\
    &=D(L(e_1)-xe_1,\ldots,L(e_n)-xe_n).
\end{align*}
For \(n\ge2\), multilinearity gives the expansion
\[
  \begin{aligned}
    p_L(x)&=D(L(e_1),\ldots,L(e_n))\\
      &\quad-x\sum_{j=1}^n
        D(L(e_1),\ldots,e_j,\ldots,L(e_n))\\
      &\quad+\cdots+(-x)^nD(e_1,\ldots,e_n).
  \end{aligned}
\]
In the sum displayed here, only the \(j\)-th input is replaced by
\(e_j\). The coefficient of \((-x)^{n-1}\) is
\[
  \sum_{j=1}^n D(e_1,\ldots,L(e_j),\ldots,e_n),
\]
where only the \(j\)-th input is acted on by \(L\).
Since \(D(e_1,\ldots,e_n)=1\), the leading coefficient is \((-1)^n\).
Thus
\[
  p_L(x)=(-x)^n+c_{n-1}(-x)^{n-1}+\cdots+c_1(-x)+c_0,
  \qquad c_0=\det L.
\]
If \(L(E)=EM\), then \((L-xI)(E)=E(M-xI)\), so
\autoref{thm:linear-algebra-i:determinant-representing-matrix} gives
\[
  p_L(x)=\det(L-xI)=\det(M-xI)=p_M(x).
\]

\subsection{Changing the basis}

Suppose \(E,F\) are bases with \(F=EH\), and write
\(L(E)=EM\), \(L(F)=FN\). Then
\[
  FN=L(F)=L(EH)=EMH=FH^{-1}MH,
\]
and hence \(N=H^{-1}MH\). For \(v=Ea=Fb\), the coordinate
columns satisfy \(a=Hb\) and \(b=H^{-1}a\).

\begin{definition}
  Two square matrices \(M,N\) are similar if there exists an invertible
  matrix \(H\) such that
  \[
    N=H^{-1}MH.
  \]
\end{definition}

\begin{proposition}
  \label{thm:linear-algebra-i:similarity-invariants}
  Similar matrices have the same determinant and characteristic polynomial.
  If \(N=H^{-1}MH\), then \(a\) is an eigenvector of \(M\) for
  \(\lambda\) if and only if \(H^{-1}a\) is an eigenvector of \(N\)
  for \(\lambda\).
\end{proposition}

\begin{proof}
  Multiplicativity gives
  \[
    \det N=(\det H^{-1})(\det M)(\det H)=\det M.
  \]
  Also \(N-xI=H^{-1}(M-xI)H\), so
  \[
    p_N(x)=(\det H^{-1})\det(M-xI)(\det H)=p_M(x).
  \]
  Finally,
  \[
    Ma=\lambda a
    \quad\Longleftrightarrow\quad
    N(H^{-1}a)=\lambda H^{-1}a.
  \]
  Since \(H\) is invertible, \(a\ne0\) if and only if \(H^{-1}a\ne0\).
\end{proof}

For \(\lambda=0\), this gives an isomorphism of kernels. Its two directions
are shown in \autoref{fig:linear-algebra-i:similar-matrix-kernels},
\[
  \ker M\longrightarrow\ker N,\quad a\longmapsto H^{-1}a,
  \qquad b\longmapsto Hb\text{ in the reverse direction}.
\]
% IMG_9322 labels the direction of its H^{-1} arrow inconsistently.
% The calculation and N=H^{-1}MH fix the directions used here.
\begin{marginfigure}
  \centering
  % Editable version of the kernel arrows in IMG_9322.
% Convention: N=H^{-1}MH. The inverse directions follow the displayed algebra.
\begin{tikzpicture}[
  x=1cm,y=1cm,
  every node/.style={font=\footnotesize},
  >={Latex[length=1.6mm,width=1.1mm]},
  line width=0.45pt
]
  \node (m) at (0,0) {\(\ker M\)};
  \node (n) at (2.6,0) {\(\ker N\)};
  \draw[->] (m.north east) to[bend left=24]
    node[above=3pt] {\(a\mapsto H^{-1}a\)} (n.north west);
  \draw[->] (n.south west) to[bend left=24]
    node[below=3pt] {\(b\mapsto Hb\)} (m.south east);
\end{tikzpicture}

  \caption[Kernel isomorphism under similarity]{For \(N=H^{-1}MH\),
    the changes of coordinates restrict to inverse isomorphisms of kernels.}
  \label{fig:linear-algebra-i:similar-matrix-kernels}
\end{marginfigure}

\begin{example}
  \label{ex:linear-algebra-i:same-polynomial-not-similar}
  The converse to characteristic-polynomial invariance fails. Let
  \[
    A=\begin{bmatrix}0&0\\0&0\end{bmatrix},\qquad
    N=\begin{bmatrix}0&1\\0&0\end{bmatrix}.
  \]
  Both characteristic polynomials are \(x^2\), but
  \[
    \ker A=\mathbb F^2,\qquad
    \ker N=\left\{\begin{bmatrix}t\\0\end{bmatrix}\;\middle|\;t\in\mathbb F\right\}.
  \]
  Their kernels have different dimensions, so they cannot be similar.
\end{example}
